\begin{lemma}[Structural equivalence is an equivalence relation]
The relation $\mathsf{PowerPresentationEq}$ is reflexive, symmetric, and
transitive on $\mathsf{PowerQ}(Q)$.
\end{lemma}
\begin{proof}
For reflexivity, unfold the recursive relation from
\noderef{PowerPresentationEq} and induct on a presentation.  An atom matches
itself by equality.  At a node, pair each child with itself and apply the
induction hypothesis in both directions, which establishes both matching
conjuncts.  For symmetry, unfold the same four tagged cases and exchange the
two matching witnesses at every node; the induction hypotheses reverse each
matched pair, and atom equality is symmetric.  For transitivity, consider
three presentations.  In the atom case, compose the equalities.  In a mixed
case, the defining relation is False.  In the node case, if a child of the
first presentation is given, the forward matching conjunct supplies a child
of the middle presentation matching it.  The forward matching conjunct for
the middle and last presentations supplies a child of the last presentation
matching that middle child.  The induction hypothesis gives a match from the
original child to the final child.  This proves the forward conjunct; the
backward conjunct is obtained by the same argument in reverse.  Hence the
recursive relation is reflexive, symmetric, and transitive.
\end{proof}
